\documentclass[11pt]{article}
\usepackage[margin=0.85in]{geometry}
\usepackage{amsmath,amssymb,booktabs,longtable,hyperref}
\usepackage[T1]{fontenc}
\hypersetup{colorlinks=true,urlcolor=blue,linkcolor=blue}
\title{Moving-Target Flow Initialization: Implementation and Closed-Loop Validation}
\author{Pan Dynamics Collision Avoidance Research}
\date{September 30, 2026}
\begin{document}
\maketitle

\section{Result and scope}
The controller now uses a moving-target fluid-inspired nonlinear rollout as a
search initialization when it has no usable prior trajectory. The seed may
collide or fail terminal admission. Its purpose is to improve the starting
point of the existing optimization, not to provide an independent safety
certificate. There remains one controller, with no driving-mode selection,
nonlinear-shooting method, or road-boundary constraints.

The new implementation completes 10 of the 14 collision-threatening fixtures.
All three previously failing 15 m/s straight-road fixtures now complete:
braking lead, crossing, and turning crossing. All completed runs have positive
independently audited sampled body clearance, with a minimum of 26.779 mm
against the unchanged 6 mm configured margin. Four circular-road fixtures still
fail at initialization because the existing permanent-cruise terminal family
cannot certify their future separation.

This is an initialization success, not a runtime improvement. The maximum
running controller call is 0.950391 s. Long avoidance and lane-recovery tails
increase repeated nonlinear validation and occasional correction work. The
50 ms control period is exceeded by 800 of 1,590 running calls. No real-time
execution claim is supported by these results.

\section{Research basis and moving-field construction}
The local source is Cheng, Li, Liu, Li and Guo, ``Virtual Fluid-Flow-Model-Based
Lane-Keeping Integrated With Collision Avoidance Control System Design,''
\emph{IEEE Transactions on Intelligent Transportation Systems}, 22(10),
6232--6241, 2021, \href{https://doi.org/10.1109/TITS.2020.2990211}{doi:10.1109/TITS.2020.2990211}.
Its PDF remains under \path{reference/}. Equations (34)--(36) concern the
stationary-obstacle field; the corresponding modified streamline has a
Gaussian lateral profile. The paper does not establish the moving-target
bicycle initialization implemented here. Its sample period is not a measured
bound on this controller's computation time.

Braquet and Bakolas, \href{https://arxiv.org/abs/2207.01747}{arXiv:2207.01747},
provide related moving-ellipse vector-field context. Their dynamics and
construction are not substituted for the project's vehicle model or claimed
as the source of the following project derivation.

Let the moving envelope be parameterized by center $c(t)$ and nonsingular
shape map $B(t)$, with $q=B^{-1}(p-c)$, $r^2=q^Tq$, and
$J=\left[\begin{smallmatrix}0&-1\\1&0\end{smallmatrix}\right]$.
For a nominal world-frame guide $v_0$, define
\begin{align}
w &= B^{-1}(v_0-\dot c-\dot Bq),\\
M(q)&=(1+r^{-2})I-2qq^Tr^{-4},\\
v_F&=\dot c+\dot Bq+B\left(M(q)w+\gamma Jq/r^2\right).
\end{align}
The translation and shape-rate terms transport the relative flow into the
world frame. Merely moving a static field's center omits these terms. Under
ideal first-order following in the smooth exterior, $h=q^Tq-1$ obeys
\[
\dot h=2q^Tw\,h/r^2.
\]
Thus the ideal boundary has zero relative normal velocity. This identity is
an algebraic check of the construction, \emph{not a required property of the
initialization trajectory}. Heading tracking, input clipping, finite
integration, and vehicle dynamics do not preserve that point-mass identity.
Inside the envelope the implementation replaces denominator $r^2$ by
$\max(1,r^2)$, keeping infeasible seed points finite without claiming safety.
The preceding derivation and ideal numerical checks were recorded separately
in ledger PDCA-20260930-004; this report records their implementation and
vehicle-level evaluation.

\section{Implementation decisions}
\begin{enumerate}
\item Retain and shift a valid prior continuation. Otherwise evaluate the
lane-feedback baseline and, when appropriate, at most two flow rollouts with
opposite tangential biases. These are bounded initialization attempts in the
same controller, not controller modes. A fully admissible candidate returns
immediately; no subsequent cost optimization is performed.
\item Predict the known target at each future hold from its fixed epoch.
Use a circular guidance envelope containing both body half-diagonals,
rectangle offsets, and the existing collision margin. The target center
velocity includes the rotational contribution of a nonzero rectangle offset.
The general field helper supports rotating/deforming ellipses; production
initialization uses a circle and zero shape-rate.
\item Activate guidance using a bounded three-second preview at
$\{0,0.5,1,1.5,2,3\}$ s. The nominal guide is cruise velocity along the given
path plus lateral recovery gain 0.5. The selected circulation magnitude is
$0.6\max(\|v_0-\dot c\|,0.5v_{\rm ref})/R$. Desired speed lies between
$\max(1.5,0.5v_{\rm ref})$ and $v_{\rm ref}$. A steering preference based on
$0.8\mu g$ and longitudinal-input preference $\pm0.35$ reduce aggressive
seed motion. These are initialization heuristics, not added optimization
constraints or a certified friction allocation.
\item Apply the actual actuator amplitude and slew limits after the soft
preferences, then integrate the existing nonlinear Fiala bicycle model.
The resulting input/state rollout supplies the dynamics, tire-force and
collision-separation linearizations together. The convex input regularizer
penalizes input deviation from this same anchor. The nominal lane/CLF
reference remains the original given path. The nonlinear ranking cost has
no residual absolute-input penalty; input proximity regularizes only the
local optimization step.
\item If a seed cannot reach the existing terminal set within 512 holds,
retain its bounded, evaluable trajectory as a restoration seed instead of
aborting before optimization. Its original terminal constraints still apply
to the final result. A cold or sufficiently changed state may request fresh
flow seeds; a raw state infinity-norm discrepancy threshold of $10^{-4}$
only controls this search choice. Smaller discrepancies still revalidate and
repair the retained plan; they are not ignored by safety admission.
\item Separate execution permission from search slack. An issued candidate
must have zero hard residual and zero prefix safety slack. The former cold
initialization slack cap was infinite and could accept a collision-containing
initialization. Positive slacks now remain restoration objects only. State
format 56 invalidates cached plans admitted under the previous semantics.
\item When a trial activates new interior collision checks, reevaluate the
incumbent with the same checks before comparing restoration scores. Restore
at least the initial trust radius for this changed check set. This removes
the stale-score rejection defect identified in the earlier diagnosis.
\end{enumerate}

The public field helper is in \path{controller/predictiveSafetyGeometry.m};
rollout and selection are in \path{controller/solvePredictiveControl.m}.
The core source count remains nine MATLAB files including configuration.
Terminal-family enlargement, bounded endpoint-polish redesign and observer
uncertainty are outside this change.

\section{Experimental method and results}
The original 14 fixtures use seven encounters at both 8 and 15 m/s. Every
given-path constant-speed baseline has a strict body overlap during the
eight-second experiment, and none starts in collision. Each completed run
executes 160 holds of 50 ms. Prefix lengths are 8 and 16 holds; maximum
completion length is 512; the call budget is 5 s and the outer iteration cap
is 24. Observations are exact, the target motion is known, road constraints
are disabled, and no random draws are used.

Plant replay uses \texttt{ode45}, relative/absolute tolerances
$10^{-11}/10^{-12}$, and 31 samples per hold. Its actual successor feeds the
next controller call. An independent Python audit recomputes target motion
and rectangle geometry at 4,960 samples per completed scenario. This is a
finite sampled validation, not a continuous-time or uncertain-plant proof.
Default slew bounds in the benchmark are infinite; a separate behavior test
uses finite steering and braking slew limits. MATLAB runs single-threaded;
timings include the complete controller call and exclude plant replay and
the independent audit. Frame zero and subsequent running calls are separate.

\begin{center}\small
\begin{tabular}{rlrrrr}
\toprule
Speed & Scenario & Old/new finish & Gap (m) & Init. (s) & Max run (s)\\
\midrule
8 & Head-on & Y/Y & 1.63940 & 1.11484 & 0.78659\\
8 & Accelerating head-on & Y/Y & 1.38929 & 0.22685 & 0.67851\\
8 & Braking lead & Y/Y & 1.13396 & 0.34902 & 0.95039\\
8 & Crossing & Y/Y & 1.15637 & 0.18703 & 0.67415\\
8 & Turning crossing & N/Y & 1.04575 & 0.28267 & 0.61383\\
8 & Curved head-on & N/N & --- & 5.01162* & ---\\
8 & Curved crossing & N/N & --- & 5.07061* & ---\\
15 & Head-on & Y/Y & 0.13959 & 1.19792 & 0.06708\\
15 & Accelerating head-on & Y/Y & 0.02678 & 0.86868 & 0.06466\\
15 & Braking lead & N/Y & 0.09355 & 0.30402 & 0.14723\\
15 & Crossing & N/Y & 1.34328 & 0.16762 & 0.06813\\
15 & Turning crossing & N/Y & 1.38045 & 0.19423 & 0.06697\\
15 & Curved head-on & N/N & --- & 5.01654* & ---\\
15 & Curved crossing & N/N & --- & 5.00166* & ---\\
\bottomrule
\end{tabular}
\end{center}
An asterisk denotes a failed initialization call, not a successful solve.
All clearance and timing columns describe the new implementation. ``Old''
is the same-environment rerun of the pre-edit snapshot at commit
\path{c217ed9518d86e7c4c678180871b6d256a4fc2dd}, whose production controller
was last changed at \path{395a8cbdc937945f6c1aef925997f731f6c9c9f8}.

The fresh old-version run completes 6/14. The previously recorded historical
run completed 7/14; its 8 m/s turning-crossing initialization was close to the
five-second budget and times out in this rerun. This timing variability is
not presented as an additional deterministic algorithmic regression. Relative
to the historical 7/14, the three clear new successes are the 15 m/s braking
lead, crossing and turning crossing.

\paragraph{Ablation.}
Disable only flow screening in a copy of the new implementation, retaining
zero-slack admission, refreshed scores and the anchor-centered input cost.
All three 15 m/s problem initializations again time out at five seconds.
This supports the contribution of the flow seed in the combined change;
it does not independently quantify each accompanying integration fix.

\paragraph{Prototype controls.}
The first 28 seed rollouts use circulation/lateral-acceleration factors
0.3/0.6. An 84-rollout screen compares parameter triples
[extra envelope radius, circulation factor, lateral-acceleration fraction]
of [0,0.6,0.8], [3.05,0.3,0.8] and [3.05,0.6,0.8]. The first triple is selected.
Extra envelope radius is not uniformly beneficial. Raw seed node clearance
is not a closed-loop safety or terminal-admission result. Both high-speed
head-on seeds can still be unsafe; their nonlinear optimization repairs them.
A further 28-rollout control uses augmented terminal feedback for the seed's
inactive recovery tail. It does not shorten the selected straight-road tails
(for example 8 m/s braking lead grows from 323 to 327 holds) and is not merged.

\section{Runtime regression and remaining failures}
The new running-call mean/median are 59.305/50.162 ms, over 1,590 calls; 800
exceed 50 ms. The old rerun's corresponding values are 16.917/6.885 ms over
954 calls, with 72 deadline misses. These aggregates cover different sets
of successful cases and must not be interpreted as a paired speed ratio.
The common 8 m/s braking-lead case nevertheless demonstrates a clear
regression: maximum running time rises from 0.117982 to 0.950391 s.

The new worst call occurs at $t=0.55$ s in that case, with a 312-hold trajectory
and one conic call. The initialization originally contained 323 holds, versus
134 for the old run. In several 8 m/s encounters flow avoidance plus return
to the small certified cruise neighborhood leaves 218--323 initial holds.
Small ODE/RK4 successor differences cause full nonlinear revalidation rather
than exact cached suffix transfer. Some frames also require endpoint
correction and a conic solve. This trace evidence explains the increased
work; no profiler decomposition assigning precise seconds to each function
was performed in this task. The historical old run had maximum 0.808422 s
in a different case that fails initialization in the fresh rerun, so neither
old maximum alone is an adequate overall performance comparison.

For the four curved fixtures, the old implementation exits before conic
optimization when its cruise endpoint cannot establish permanent separation.
The new code permits uncertified seeds and enters restoration, but still
reaches the soft time limit without a witness. The unchanged circular
terminal certificate separates complete future orbit envelopes and is
independent of phase in these branches. Its negative margins, approximately
$-3.5436$ and $-39.748$ m, are certificate lower bounds, not observed penetration
depths. Flow initialization cannot repair a negative constant terminal row.
For opposing travel around the same circle, perpetual nominal cruising also
truly meets again eventually; timing-aware geometry alone does not make that
particular continuation safe. These findings and their phase checks are
documented in \path{report/REMAINING_FAILURE_DIAGNOSIS_20260930.md}.

\section{Validation, evidence and limitations}
The final rerun after removing the residual absolute-input ranking term
preserves the ten successes and their sampled minimum clearances. All timing
figures above come from this final run. All 273 selected tests pass across nine classes. New coverage checks moving
boundary transport, world-velocity shifts, rotational equivariance, finite
interior guidance, six known-target admissions, finite actuator slew, unsafe
flow seeds repaired by optimization, and rejection of positive-slack
execution. Source-budget and existing vehicle, terminal, geometry and
scenario tests are included. Code analysis has only the existing 15 sparse
indexing performance advisories in the solver among the five scanned files;
it reports none in the other four.

An initial one-frame threat harness was rejected because its 50 ms baseline
did not contain a collision. An early implementation run exposed an empty
struct-array insertion error in seed diagnostics; it was fixed before the
reported campaign. The first selected test run passed 272/273, with one
legacy test still expecting positive-slack execution; the updated behavior
test and all other selected tests pass on rerun. Failed attempts remain
recorded rather than counted as successful simulations.

Compact evidence is in \path{report/MOVING_FLOW_INITIALIZATION_20260930/}:
campaign and ablation summary, both independent audits, per-call timing CSV,
fixture configurations, tests, analyzer results, prototype controls, source
hashes, raw-output hashes and reproduction instructions. Full poses, MAT
snapshots and copied/prototype sources remain in the named external cache.
No solver dependency, generated binary, unrelated estimator change or
manuscript artifact belongs to this implementation commit.

The accepted scope is a useful initialization heuristic with demonstrable
closed-loop gains in the three high-speed straight-road failures. It does
not solve the four terminal-family failures, certify between-sample or
uncertain-plant safety, or meet the 50 ms runtime target.
\end{document}
